The DyninstAPI library provides an interface for instrumenting and working with
binaries and processes. The user writes a mutator, which uses the DyninstAPI
library to operate on the application. The process that contains the mutator
and DyninstAPI library is known as the mutator process. The mutator process
operates on other processes or on-disk binaries, which are known as mutatees.

The API is based on abstractions of a program. For dynamic instrumentation, it
can be based on the state while in execution. The two primary abstractions in
the API are points and snippets. A point is a location in a program where
instrumentation can be inserted. A snippet is a representation of some
executable code to be inserted into a program at a point. For example, if
we wished to record the number of times a procedure was invoked, the point
would be entry point of the procedure, and the snippets would be a statement
to increment a counter. Snippets can include conditionals and function calls.

Mutatees are represented using an address space abstraction. For dynamic
instrumentation, the address space represents a process and includes any
dynamic libraries loaded with the process. For static instrumentation, the
address space includes a disk executable and includes any dynamic library
files on which the executable depends. The address space abstraction is
extended by process and binary abstractions for dynamic and static
instrumentation. The process abstraction represents information about a
running process such as threads or stack state. The binary abstraction
represents information about a binary found on disk.

The code and data represented by an address space is broken up into function
and variable abstractions. Functions contain points, which specify locations
to insert instrumentation. Functions also contain a control flow graph
abstraction, which contains information about basic blocks, edges, loops, and
instructions. If the mutatee contains debug information, DyninstAPI will also
provide abstractions about variable and function types, local variables,
function parameters, and source code line information. The collection of
functions and variables in a mutatee is represented as an image.

The API includes a simple type system based on structural equivalence. If
mutatee programs have been compiled with debugging symbols and the symbols
are in a format that Dyninst understands, type checking is performed on code
to be inserted into the mutatee. See Section for a complete description of the
type system.

Due to language constructs or compiler optimizations, it may be possible for
multiple functions to overlap (that is, share part of the same function body) or
for a single function to have multiple entry points. In practice, it is
impossible to determine the difference between multiple overlapping functions
and a single function with multiple entry points. The DyninstAPI uses a model
where each function has a single entry point, and multiple functions may overlap
(share code). We guarantee that instrumentation inserted in a particular function
is only executed in the context of that function, even if instrumentation is
inserted into a location that exists in multiple functions.
